\documentclass[runningheads]{llncs}
\usepackage[T1]{fontenc}
\usepackage{lmodern}
\usepackage{amsmath,amssymb,booktabs,array,graphicx}
\usepackage{url}
\usepackage[hidelinks]{hyperref}
\newcolumntype{L}[1]{>{\raggedright\arraybackslash}p{#1}}
\newcommand{\method}{ESV-Gap}
\begin{document}
\title{ESV-Gap: Scope-Aware Decision Control for Autonomous Scientific Agents}
\titlerunning{Scope-Aware Evidence Control for Scientific Agents}
\author{Anonymous Author(s)}
\authorrunning{Anonymous Author(s)}
\institute{Anonymous Institution}
\maketitle

\begin{abstract}
An autonomous scientific agent must decide when to suppress a proposed empirical question and when to seek further evidence. Topical correspondence and measurements from different experiments can incorrectly trigger suppression. ESV-Gap supplies an explicit scope contract: one reviewed experiment must jointly meet the task, domain, condition, and outcome obligations before closing a question. Otherwise the controller retains residual obligations and abstains when review or the declared search procedure is incomplete. We establish conditional witness integrity, noncomposition, restricted closure persistence, and conservative handling of incomplete searches, checking 8,192 constructed two-record inputs and additional metamorphic cases. A source-linked grounding retriever supports the review stage. On 42 expert-labelled Context24 claims from 31 papers, leave-one-paper-out ROUGE-L F1 is 0.2122 versus 0.1805 for BM25 under the same 512-word budget; the descriptive paired-paper difference is 0.0393 with a 95\% interval [0.0131, 0.0666]. Ablations and cross-lab transfer expose limitations of the learned context prior. A separate three-domain retrospective failure analysis motivates the controller without serving as its evaluation. Real-literature closure accuracy and end-to-end agent utility remain unmeasured; the control guarantees depend on trusted annotations and do not certify global novelty.
\keywords{Scientific agents \and Evidence scope \and Hypothesis triage \and Abstention \and Literature-based discovery}
\end{abstract}

\section{Introduction}
Literature-based discovery seeks connections between published findings~\cite{swanson1986}. Scientific AI assistants now generate questions, methods, and experimental plans using scholarly corpora and academic graphs~\cite{baek2025}. Their autonomy creates a practical decision problem: which proposed questions warrant further formulation, which have already been answered within a specified scope, and which require more evidence? A plausible candidate is only the beginning of that decision.

A missing graph edge can reflect an unperformed experiment, an incomplete abstract, an extraction error, or inconsistent terminology. Likewise, a later paper mentioning a candidate's terms may discuss an unresolved limitation, propose a conceptual framework, or evaluate a different task. Finding an action verb such as ``addresses'' does not establish that the paper achieved a claimed outcome under the required conditions. Even a genuine measurement may be irrelevant to an underspecified question.

The original \method{} architecture uses relation extraction, candidate generation, external checking, and an abstention policy. Its small IoT case study motivated a broader audit, but did not establish end-to-end discovery efficacy. In the preserved primary run, 133 raw signals were routed to zero automatically eligible candidates, five review candidates, and 128 rejected signals. These counts describe the stored gate output, not the five-case semantic reinterpretation used in the earlier manuscript. Keeping those stages separate avoids presenting a selected case proportion as population performance.

We investigate three questions. \textbf{RQ1:} Which explicit evidence obligations and control properties govern a scoped suppression decision? \textbf{RQ2:} Does methodological-context information improve retrieval of expert grounding quotations under a common budget? \textbf{RQ3:} What failure modes of earlier candidate generation and proxy matching motivate this control boundary? We contribute (i) a scope controller with conditional properties and executable counterexample checks, (ii) a source-linked context retriever with a paper-held-out public-benchmark evaluation and ablations, and (iii) a regenerated three-domain diagnostic audit and a frozen review pilot. The retrieval comparison measures an auxiliary acquisition task. It does not evaluate numerical witness extraction, autonomous closure, or end-to-end research-gap discovery. The system makes the automatic decision conditional on semantic review; it has not demonstrated fully autonomous evidence understanding.

\section{Related Work and Contribution Boundary}
ResearchAgent combines language models, academic-graph information, and iterative reviewing agents to develop ideas~\cite{baek2025}. It includes human and model-based evaluation; we therefore avoid the blanket claim that prior scientific agents lack verification. Our narrower concern is the explicit evidence condition that permits a controller to suppress or retain a particular empirical question.

M\"uller-Bloch and Kranz's research-gap framework already includes localization, characterization, verification, and presentation~\cite{mullerbloch2015}. Its six types include contradictory evidence, knowledge void, action--knowledge conflict, methodological conflict, evaluation void, and theory application void. Our scope contract addresses empirical capability and evaluation questions; it is not a complete operationalization of that typology. Theoretical and methodological questions may require proofs or arguments rather than numerical witnesses.

SciFact studies scientific claim verification~\cite{wadden2020}; SciFact-Open studies retrieval in an open domain where evidence may support only a special case~\cite{wadden2022}. Context24 distinguishes stating a claim from empirically grounding it in methods, figures, or tables~\cite{chan2024}. We reuse its expert methodological quotations for an auxiliary retrieval evaluation, with a distinct scoring protocol rather than an official shared-task result. Our lexical prior and redundancy penalty are simple retrieval components; their combination is not claimed as a new general retrieval algorithm. The contribution is their measured role and limits within a source-linked, scope-controlled workflow. A source substring verifies where text came from, not whether it entails an annotated measurement or condition.

BM25 supplies a transparent lexical comparator~\cite{robertson2009}. Scientific document representations such as SPECTER and evaluation suites such as SciRepEval provide alternatives for a subsequent dense-retrieval comparison~\cite{cohan2020,singh2023}; they have not been run here. Graph embeddings and communities can prioritize missing relationships~\cite{bordes2013,blondel2008}, but graph structure alone cannot establish experimental closure. Selective prediction motivates abstention~\cite{geifman2017}; our deterministic rule supplies neither a calibrated probability nor a population risk guarantee.

\section{Grounding Acquisition and Scoped Evidence Control}
Figure~\ref{fig:system} distinguishes the agent's proposed question, evidence acquisition, trusted review, and automatic action. The implemented controller consumes reviewed experiment records and a search ledger; neither graph structure nor retrieval scores authorize suppression.

\begin{figure}[t]
\centering
\setlength{\unitlength}{1mm}
\begin{picture}(112,85)
\scriptsize
\put(0,70){\framebox(34,14){\shortstack{Agent-proposed\\empirical question}}}
\put(39,70){\framebox(34,14){\shortstack{Reviewed scope\\contract: $H$}}}
\put(78,70){\framebox(34,14){\shortstack{Evidence acquisition:\\documents and spans}}}
\put(34,77){\vector(1,0){5}}
\put(73,77){\vector(1,0){5}}
\put(95,70){\vector(0,-1){13}}
\put(78,43){\framebox(34,14){\shortstack{Provenance and\\semantic review}}}
\put(39,43){\framebox(34,14){\shortstack{Single-experiment\\witness check}}}
\put(0,43){\framebox(34,14){\shortstack{Search ledger:\\required probes}}}
\put(56,70){\vector(0,-1){13}}
\put(34,50){\vector(1,0){5}}
\put(78,50){\vector(-1,0){5}}
\put(56,43){\vector(0,-1){8}}
\put(39,24){\framebox(34,11){\shortstack{Disposition and\\residual obligations}}}
\put(56,24){\vector(0,-1){8}}
\put(0,2){\framebox(112,14){\shortstack{Complete witness: suppress within scope\\Otherwise: request residual evidence or abstain for review}}}
\end{picture}
\caption{Control boundary for scientific-agent hypothesis triage. Evidence-document retrieval and within-paper span proposal are distinct acquisition steps; Context24 evaluates only the latter. Semantic review is a trusted input. The diagram represents the workflow, not a measured end-to-end autonomous loop.}
\label{fig:system}
\end{figure}

\subsection{Source-Linked Methodological Context}
Given a claim and its source paper, we propose passages for semantic review before constructing experiment records. This within-paper step assumes the source document is available; it does not discover witness papers or decide whether a research gap is open. Split the unmodified text into windows of at most 96 whitespace words at stride 80, retaining original character offsets. Let $B(w,q)$ be BM25 claim relevance and let $c_+(v)$ and $c_0(v)$ count token $v$ in training methodological quotations and training full texts, respectively. Training excludes every claim and the full text belonging to the held-out paper. With vocabulary $V$ and totals $N_+$ and $N_0$, define
\begin{equation}
 \ell(v)=\max\!\left(0,\log\frac{(c_+(v)+1)/(N_++|V|)}
 {(c_0(v)+1)/(N_0+|V|)}\right),\quad
 A(w)=\frac{\sum_{v\in w}\ell(v)}{\max(1,|w|)}.
\end{equation}
Tokens occurring fewer than twice in training quotations receive zero weight. Tokenization uses lowercased, Unicode-normalized ASCII alphanumeric tokens. Normalize $B$ and $A$ separately by their maximum across the source windows, using zero when that maximum is zero. Given selected windows $S$, greedily select a nonoverlapping window maximizing
\begin{equation}
 F(w\mid q,S)=0.5\widetilde B(w,q)+0.5\widetilde A(w)
 -0.2\max_{s\in S}J(w,s),
\label{eq:retrieval}
\end{equation}
where $J$ is token-set Jaccard similarity and the empty-set penalty is zero. Ties prefer earlier source offsets. Stop at the word budget, truncating the last selected quotation if necessary; present selected spans in source order. Weights were fixed in a local protocol before the comparative run, without held-out score tuning. This exploratory protocol was not preregistered.

Returned spans retain \texttt{reviewed=false}. They are acquisition proposals for a reviewer to map to tasks, conditions, measurements, and experiment membership. A useful method quotation may omit every numeric outcome needed for closure. The retrieval module therefore cannot supply a complete witness merely by finding relevant text; the controller below separately requires reviewed annotations.

\subsection{Predeclared Empirical Question}
Represent a question by
\begin{equation}
 H=(\tau,d,C,R,t),\quad R=\{(m_j,\circ_j,b_j,u_j)\}_{j=1}^{J},
\end{equation}
where $\tau$ is the task, $d$ the domain, $C$ a nonempty set of required experimental conditions, $R$ a nonempty set of outcome requirements, and $t$ the evidence-availability cutoff. Each requirement names a metric $m_j$, a comparison $\circ_j\in\{\leq,\geq\}$, a finite threshold $b_j$, and a unit $u_j$. Requirements must be reviewed and fixed before assessing candidate witnesses. An underspecified phrase such as ``Retrofitting--Security'' fails this contract; the controller requests clarification rather than inventing a task or a success threshold.

For illustration only, a question might ask whether an IoT detector identifies unknown attacks on edge hardware while achieving F1 at least 0.90 and latency at most 10\,ms. These are synthetic test thresholds, not validated deployment requirements or measured outcomes from our corpora. If latency is reported on a workstation, the hardware condition is unfulfilled. If one experiment reports F1 and another reports latency, they cannot be combined into a complete witness.

\subsection{Reviewed Witness and Residual Obligations}
An evidence record has a paper identifier, experiment identifier, recorded availability date, source text, and annotations for task, domain, conditions, and measurements. Each asserted facet carries an exact source span and a semantic-review flag. Metadata and experiment membership also require review: membership is the assertion that all annotated measurements concern the same experiment. The current implementation trusts those reviewed annotations; it does not automatically infer entailment or experiment identity from text.

Let $P(e;E,H)$ mean that identifiers, availability metadata, experiment membership, and all required quoted annotations pass these checks, and that $e$'s paper--experiment key occurs exactly once in the supplied multiset $E$. Let $s_f(e,H)$ denote a satisfied facet. A within-scope witness is
\begin{equation}
 W(e,H;E)=P(e;E,H)\land a(e)\leq t\land s_{\tau}(e,H)\land s_d(e,H)
 \land\bigwedge_{c\in C}s_c(e,H)\land\bigwedge_{r\in R}s_r(e,H).
\label{eq:witness}
\end{equation}
The quantifier is $\exists e\in E\,W(e,H;E)$: all obligations hold in one supplied experiment record. It is not $\bigwedge_f\exists e_f\in E\,s_f(e_f,H)$, which could silently combine different studies, datasets, or hardware configurations. Duplicate records for the same paper--experiment key are flagged for review rather than merged. The dependence on $E$ matters for the insertion property below.

For every non-witness, the audit records residual states: unreported, unreviewed or unquoted, different scope, incomparable units, or a failed threshold. Values with incompatible units are retained as unresolved rather than silently converted. Missing outcome data do not imply that a capability is impossible or scientifically novel. An annotation with a valid quotation but no semantic review is also unresolved.

\subsection{Bounded Search and Action Policy}
A search ledger declares required probe identifiers before recording results. Each required probe must have a provider, query, snapshot identifier, completed status, completed pagination, and completed semantic review. The ledger cutoff must equal the contract cutoff. Completion means the declared procedure completed; it is not exhaustive coverage of the literature, nor proof that a modern index reproduces a historical one.

Table~\ref{tab:policy} gives the implemented precedence. A complete witness can suppress an empirical absence candidate even if another probe failed: one valid existential witness suffices to close the specified question. Without such a witness, an incomplete search or unchecked evidence forces review. Reviewed evidence that satisfies only part of the contract produces residual obligations. A completed search with no witness yields only an open review queue.

\begin{table}[t]
\caption{Scope-controller policy. Every branch records \texttt{novelty\_established=false}. Closure is relative to the reviewed contract and annotations.}
\label{tab:policy}
\centering\small
\begin{tabular}{L{4.4cm}L{3.1cm}L{3.5cm}}
\toprule
Condition, in precedence order & Disposition & Action \\
\midrule
Invalid or unreviewed contract & Review required & Abstain and clarify \\
At least one complete witness & Closed within scope & Suppress candidate \\
Incomplete protocol or unchecked evidence & Review required & Abstain and escalate \\
Reviewed partial evidence & Partial evidence & Request residual evidence \\
No witness after completed protocol & Open for review & Abstain and escalate \\
\bottomrule
\end{tabular}
\end{table}

\begin{proposition}[Conditional witness integrity]
For a valid contract, the implemented controller returns \textnormal{closed within scope} only if one supplied record satisfies every obligation in Eq.~\ref{eq:witness}, relative to the trusted semantic and metadata annotations.
\end{proposition}
\begin{proof}
The controller checks provenance and availability before evaluating a record. It evaluates every required facet and measurement in that record and marks it as a witness only when every state is satisfied. The closure branch requires at least one such witness. No branch merges record facets. Thus closure implies a complete individual witness. Semantic truth of the supplied annotations is an assumption, not a conclusion of this argument.
\end{proof}
This is a property of the control rule. It neither establishes that source annotations are correct nor proves novelty when no witness is found. Vocabulary saturation, repeated limitation statements, and language-model confidence do not enter the witness rule as substitutes for empirical evidence.

\subsection{Conditional Control Properties}
Write $D(H,E,L)$ for the disposition under contract $H$, evidence multiset $E$, and ledger $L$. All properties assume structurally valid inputs and trusted annotations. For a reviewed valid contract, closure is equivalent to $\exists e\in E\,W(e,H;E)$; the implementation gives that branch precedence over ledger incompleteness.

\begin{proposition}[Noncomposition and conservative incomplete search]
If no individual record is a complete witness, the controller cannot close the question even if the union of satisfied facets covers every obligation. If additionally the ledger is incomplete, it returns \textnormal{review required} and abstains.
\end{proposition}
\begin{proof}
The witness predicate conjuncts obligations within each record, so their union across records never satisfies the existential closure test. With that branch unavailable, ledger incompleteness selects the review branch before any partial or open disposition.
\end{proof}

\begin{proposition}[Restricted closure persistence and contract strengthening]
For fixed $H$, a closure persists after adding records if at least one original witness is unchanged and its key remains unique. For fixed $E,L$, let $H'$ retain task, domain, cutoff, and units, add required conditions or metrics, and only tighten existing thresholds. If both contracts are reviewed and valid, closure under $H'$ implies closure under $H$.
\end{proposition}
\begin{proof}
In the first case, the preserved witness retains its provenance and all satisfied obligations; its unchanged unique key preserves $P$. In the second, every witness satisfying the additional conditions and tighter bounds satisfies the original obligations, which are a subset or weaker. It therefore triggers the original closure branch.
\end{proof}
Unrestricted insertion monotonicity is false: adding a duplicate of the sole witness invalidates its unique-key check and can require review. Review outcomes are also not monotone, because a new valid witness may close a previously incomplete case. These exceptions are deliberate consequences of the policy, not annotation-accuracy guarantees.

The rule is deterministic for an identical serialized input. Re-evaluating it yields the same result and audit hash; permuting evidence records preserves disposition and action but may reorder traces and change the hash. This is repeatability of a decision function, not idempotence of an autonomous agent acting on its output. No claim is made that repeated search or contract revision preserves the decision.

\section{Expert-Labelled Grounding Retrieval Evaluation}
\subsection{Dataset, Split, and Scoring}
We use all 42 labelled Task-2 example claims from Context24, covering 31 source-paper identifiers: 28 claims from Akamatsu Lab and 14 from MegaCog Lab. Full-text parses are available for every example. For each held paper, learn the prior from only the other papers and evaluate all its claims. The downloaded 109-claim official test file omits gold contexts and is neither scored nor used for tuning. This is cross-validation on the small public labelled example set, not a new independent test collection.

Every arm receives identical source windows and at most 512 whitespace words. Comparators are BM25 ($k_1=1.2$, $b=0.75$), sublinear TF-IDF cosine, and prior-only ranking. Ablations add diversification to BM25 or remove it from Eq.~\ref{eq:retrieval}. No language-model-generated labels are used. We concatenate predictions in source order and gold quotations in stored order, computing non-stemmed word ROUGE-1, ROUGE-2, and ROUGE-L F1. These lexical measures are sensitive to wording and quotation order. The official scorer instead matches each predicted snippet to its best gold snippet, uses stemmed ROUGE, and also computes BERTScore; our numbers are not comparable to the leaderboard.

Table~\ref{tab:context} reports claim-macro means. For uncertainty, average paired ROUGE-L differences within each paper, then resample the 31 paper means with replacement for 2,000 draws (seed 20261001). This paper-macro estimand differs from the claim-macro table difference. The percentile interval is descriptive: overlapping cross-validation training sets, two laboratories, and a small convenience sample preclude a population performance guarantee.

\begin{table}[t]
\caption{Context24 labelled-example retrieval, 42 claims/31 papers. All methods use 512 words and the same windows. Scores are concatenated, non-stemmed word F1, not official shared-task scores.}
\label{tab:context}
\centering\small
\begin{tabular}{lrrr}
\toprule
Method & ROUGE-1 & ROUGE-2 & ROUGE-L \\
\midrule
BM25 & .3573 & .1455 & .1805 \\
TF-IDF cosine & .3639 & .1445 & .1857 \\
Context prior only & .4065 & .1976 & .2045 \\
BM25 + diversification & .3575 & .1436 & .1805 \\
Context hybrid, no diversification & .4125 & .1963 & .2102 \\
Context hybrid + diversification & .4132 & .1991 & .2122 \\
\bottomrule
\end{tabular}
\end{table}

\subsection{Ablations, Budget Sensitivity, and Transfer}
The proposed retriever improves claim-macro ROUGE-L over BM25 by 0.0317; its paper-macro difference is 0.0393, with a 95\% percentile interval [0.0131, 0.0666]. The context prior provides most of the gain. Post-hoc paired-paper comparisons against prior-only and the nondiversified hybrid give differences of 0.0019 and $-0.0005$, with intervals [$-0.0274$, 0.0337] and [$-0.0060$, 0.0044]. These diagnostics provide no clear incremental advantage from claim mixing or diversification over those ablations. We retain the frozen method and report this ambiguity rather than selecting a new weight after evaluation.

Within-lab ROUGE-L for the proposed method versus BM25 is 0.2297/0.1850 for Akamatsu and 0.1771/0.1714 for MegaCog. TF-IDF reaches 0.1918 on MegaCog and exceeds the proposed method there. In post-hoc budget checks, proposed/BM25 scores are 0.1878/0.1685 at 256 words and 0.2144/0.1779 at 768; neither budget changes the original 512-word configuration.

To examine reliance on domain vocabulary, a further post-hoc diagnostic trains only on the other laboratory, excluding all held-lab paper identifiers. Transfer to Akamatsu (14 training claims, six training papers) gives 0.1873 versus BM25's 0.1850. Transfer to MegaCog (28 training claims, 25 training papers) gives 0.1504 versus 0.1714. Thus the gain does not transfer consistently across these groups. In an unfamiliar domain the learned prior needs separate validation; BM25 remains an available acquisition baseline. These results do not justify deploying the prior across the audit domains without evaluation.

All 252 saved predictions pass exact source-span and word-budget checks. A separate verification script recounts aggregate arithmetic; ten bounded token-prefix checks compare the optimized ROUGE-L implementation with conventional dynamic-programming LCS. Six retrieval unit tests check offsets, budgets, scoring edge cases, and empty inputs. These implementation checks do not establish semantic sufficiency or downstream closure accuracy.

\section{Retrospective Failure Analysis}
\subsection{Primary Diagnostic Scope and Protocol}
This analysis explains why proxy matching cannot authorize suppression; it evaluates the earlier retriever, not the new controller. The preserved four-domain archive contains 378 keyed records and 3,188 triples. One matched MongoDB record has unresolved source identity and publication year, so our primary diagnostic analysis excludes the entire 53-record MongoDB domain. This exclusion was decided after inspecting archived results. It is not a preregistered sampling decision, provenance repair, or certification that the remaining domains are clean.

We rerun candidate generation, silver-control construction, and confidence/BM25 rankings on 325 records and 2,545 triples: 150 IoT intrusion-detection, 143 microservice-security, and 32 handwritten-mathematical-expression papers. Cutoffs are 2022, 2023, and 2024, requiring at least ten records on each side. Eight folds qualify; seven have positive controls. Regenerated candidates, controls, and metrics equal archived outputs for every retained fold. Input, configuration, and code hashes are checked; historical files remain unchanged.

Pre-cutoff triples generate candidates; future triples generate silver controls through limitation-plus-action correspondence or newly instantiated typed relations. These are not expert-confirmed research questions or resolutions. Publication years are stored metadata, not reconstructed first-online dates. Support-ID checks exclude future candidate sources but do not reconstruct historical model training or search indexes.

Comparators rank direct \texttt{LACKS} triples by extractor confidence or BM25 over subject/capability text ($k_1=1.2$, $b=0.75$). BM25's domain-label query plus ``limitation unresolved challenge'' was chosen post hoc. It reranks the extracted pool and is not an independent raw-text discovery baseline. Matching uses capability token similarity at least 0.62 and subject similarity at least 0.25, or capability similarity at least 0.82; similarity combines containment and Jaccard overlap. Proxy recall counts matched positive controls under the same candidate budget, averaged across the seven scored folds.

\subsection{Diagnostic Findings and Excluded Secondary Results}
\begin{table}[t]
\caption{Primary three-domain diagnostic audit, seven scored folds. Controls are silver and non-independent. Scores concern the earlier candidate retriever, not grounding retrieval or scope-controller accuracy.}
\label{tab:macro}
\centering\small
\begin{tabular}{lrrrr}
\toprule
Ranking & R@5 & R@10 & R@20 & Matches@20 \\
\midrule
Earlier ESV-Gap & 0 & 0 & .0179 & 1/61 \\
Confidence--LACKS & 0 & 0 & 0 & 0/61 \\
BM25--LACKS & 0 & .0179 & .0485 & 4/61 \\
\bottomrule
\end{tabular}
\end{table}

The 61 control rows contain 29 distinct control IDs within domain; repeated cutoffs share papers. The retriever produces 38 candidate rows, all explicit limitations, while its typed evidence-map branch yields zero. The direct-limitation pool contains 396 rows across folds. Three folds yield no ESV-Gap candidates. Table~\ref{tab:macro} shows weak candidate utility under the proxy; it supplies no positive evidence of graph-specific discovery value.

The four BM25 matches are examined at abstract level with AI assistance. The IoT adversarial-attack match concerns a cybersecurity synthesis~\cite{jayawardena2026}; a microservice match concerns a systematic mapping study~\cite{rahaman2023}; and an API/service-mesh match concerns a Kubernetes overview~\cite{kampa2024}. None of those abstracts reports a matched same-scope resolution experiment. The remaining match reports quantitative DoS detection~\cite{olaya2024}, but its candidate, ``Retrofitting--Security,'' lacks that task and required conditions. These observations are not independent labels or four verified false positives. Full texts may contain evidence absent from abstracts.

For transparent secondary accounting, the original four-domain archive retains ten executable folds, nine scored folds, 64 control rows, and 39 ESV-Gap candidate rows. Its macro R20 is 0.0139 for ESV-Gap and 0.0933 for BM25--LACKS, with five BM25 matches; the fifth is the disputed MongoDB record. These values are excluded from the primary table. The three-domain score ordering persists without that domain. Exclusion addresses its influence on the primary comparison, while its identity/date and all silver labels still need verification. The regenerated three-domain analysis supersedes the earlier aggregate-only omission check.

\section{Implemented Extension and Verification Status}
\subsection{Controller Checks}
The new Python controller consumes a JSON contract, evidence records, and a search ledger, producing dispositions, residual traces, and an input SHA-256. It is an explicit opt-in component; it does not retroactively replace the earlier decision engine or alter historical runs. It accepts no vocabulary-saturation rule as evidence of experimental completeness. Boolean values are excluded from numeric measurements, nonfinite values never satisfy requirements, and incompatible units remain unresolved.

Twenty-five constructed tests pass. They check a complete individual witness, failure to combine measurements across experiments or papers, future-record exclusion, missing dates, unreviewed metadata and experiment membership, fabricated or unchecked quotes, wrong task/domain, missing or failed conditions, incompatible units, nonfinite values, failed outcome thresholds, underspecified or unreviewed contracts, missing probes, incomplete pagination, an empty completed search, witness precedence despite a failed probe, input immutability, duplicate experiments, and malformed metric names. These are software counterexamples and rule checks. Their pass rate is not an estimate of evidence-extraction accuracy, real-world hallucination suppression, or scientific-discovery success.

A separate bounded verifier enumerates 64 symbolic states per record: task match, domain match, condition satisfaction, two outcome thresholds, and experiment-review status. Two records with distinct keys and two ledger-completion states produce $64^2\times2=8,192$ inputs. It checks closure against the known per-record conjunction and incomplete-search abstention in every case. Among the two-record state pairs, 180 have reviewed records whose union covers all five obligations but neither record is a witness; all remain nonclosed under both ledger states. Additional checks cover 64 unique-key insertions, 64 record permutations, 64 repeat evaluations, 144 tightened-threshold combinations, and the duplicate-key counterexample. Dates, units, quoted provenance, and metadata are fixed valid values in this enumeration; the unit tests vary these separately. The enumeration is finite constructed verification, not full state-space model checking or real-literature effectiveness.

\subsection{Extraction-Free Evidence Retrieval and Blinded Review Preparation}
For each of the 39 archived candidate rows, a new local retriever forms a query from the pre-cutoff subject and capability and ranks raw titles plus abstracts using BM25. It excludes the candidate's supporting paper identifiers and all records with stored years later than the cutoff; up to five nonzero-score records are retained. This produces 188 candidate--evidence pairs involving 97 distinct evidence-paper identifiers. They are preparation and workload counts, not efficacy metrics. Evidence ranking uses no extracted triples; candidate creation still does.

Two independently randomized CSV packets omit ranking scores and system judgments. Six annotation labels distinguish same-scope closure, partial evidence, supporting limitation, topical-only correspondence, insufficient evidence, and an underspecified candidate. Reviewers must supply quoted evidence and an experiment identifier for a closure label. All fields are initially blank; 39 contract templates deliberately leave task, conditions, and numerical requirements for review rather than generating unsupported success criteria. An agreement scorer rejects missing labels and overlapping reviewer identities, computes nominal Cohen's $\kappa$, and lists disagreements. It does not mistake agreement for scientific correctness or system accuracy.

For a manageable pilot, we freeze a subset before obtaining any labels. Excluding MongoDB, group candidates by normalized subject/capability within domain and retain their earliest cutoff; then select four per domain with a fixed seed, keeping all prepared evidence pairs. This yields 12 candidate rows, 53 pairs, and 40 evidence-paper IDs. Two separately shuffled packets hide scores and judgments. The selection uses no future controls; it does not pool proposals from independent discovery systems or justify statistical power. Reviewer contracts must be fixed from pre-cutoff sources before assessing evidence.

Independent assessment has not occurred. Evidence is ranked within retained corpora, and origins are excluded by identifier rather than author group or complete preprint equivalence. Neither the 188-pair preparation nor the 53-pair pilot establishes external corroboration, comprehensive search, or a positive controller result.

\section{Remaining Closure Evaluation}
A prospective study should freeze scoped questions and requirements using only pre-cutoff information, then independently review evidence with system identity hidden. Questions should be pooled from direct limitations, an extraction-free proposal baseline, and graph proposals under equal candidate budgets. The same witness rules and retrieval budgets must apply to each arm. Dense and hybrid evidence retrievers can be added with frozen versions and query policies. A lexically strong ranking should not receive credit for a closure that lacks the required experimental scope.

At least two domain-qualified reviewers should first assess task/condition specificity, then annotate source-linked facets and experiment membership. Disagreements need adjudication by a third reviewer or a documented consensus process. Full text should be inspected for potential witnesses and abstracts marked insufficient when needed. The frozen 53-pair subset is a pilot instrument, not a sample size justified by a desired precision interval. Cases share candidates and papers; statistical analysis should account for that grouping.

The required endpoints are witness-document recall within a declared corpus, adjudicated closure precision, decision coverage over valid scoped questions, expert-assessed proposal utility, and retrieval plus review cost. Closure precision is undefined if no closures occur. None of these controller endpoints is estimated by the Context24 lexical scores. The grounding experiment assumes the source paper is already given, whereas witness-document retrieval must first find relevant papers.

Ablations should remove one requirement at a time: membership in one experiment, checked metadata, semantic review, or completion of the search. Studies under different conditions, reviews, conceptual proposals, and already resolved limitations should be explicit hard negatives. Retain graph complexity only if it adds expert-confirmed candidate utility. Keep open cases unresolved.

\section{Limitations and Reproducibility}
The primary audit uses three retained corpora of unequal size and noisy silver controls; the fourth domain is excluded post hoc because of unresolved provenance. Recorded years remain a coarse historical proxy. Repeated folds reuse papers; no independence-based confidence intervals or significance tests are justified. The BM25 limitation comparator and domain exclusion were post hoc. The new raw-abstract evidence retriever is not an end-to-end gap-generation comparator and has no human relevance labels.

The Context24 experiment uses only 42 public labelled examples, two laboratories, and source papers supplied by the dataset. Its cross-validation prior can exploit methodological vocabulary shared with the other papers; the cross-lab failure demonstrates a concrete transfer risk. Citation overlap, duplicate publication versions, and author dependence are not independently audited. We measure lexical overlap with selected expert quotations, which are not exhaustive witness annotations. Dense retrieval and official Context24 semantic scoring are not evaluated. The retrieval gain cannot be interpreted as a reduction in false closure decisions.

The scope controller depends on correct semantic, numerical, and metadata annotations. Exact quotations cannot prove entailment, and a reviewed flag does not independently verify its reviewer. A numeric contract is appropriate only when domain experts can defend meaningful metrics and thresholds. It may miss qualitative discoveries or theoretical contributions. Our conditional witness-integrity argument and synthetic tests supply control-rule evidence; they do not estimate real-literature closure accuracy or generalization.

Retained reports contain SHA-256 fingerprints of corpus/triple inputs, configuration, and replay code. Preparation reports additionally hash the four corpora and ten archived fold reports. Table~\ref{tab:macro} is regenerated from the retained inputs, and its local report includes per-fold metrics and code/environment hashes; reproducing it requires those inputs and dependencies. The public Context24 data are downloaded at HuggingFace revision \texttt{457d3b5cb4bb8ade34e37458f4900c6eae0959bb}, licensed CC BY 4.0~\cite{contextdata}. Its report records input/code/protocol hashes, per-claim scores, source spans, and the Python version. The retrieval implementation uses only the Python standard library. Scripts, reports, packets, and source are retained locally. No public archival release of our artifact, independent reproduction, redistributability of the four audit corpora, or historically complete search snapshot is claimed.

\section{Conclusion}
ESV-Gap separates acquisition of empirical context from a scoped evidence decision. On a small public expert-labelled benchmark, the context prior improves lexical grounding retrieval over BM25 at a matched word budget; ablations and cross-lab transfer reveal limited evidence for diversification and domain portability. The controller requires one reviewed experiment to meet every declared obligation and keeps absence claims in review. The primary three-domain failure analysis exposes weak earlier candidate availability and unreliable closure proxies. Conditional control properties, constructed checks, and a frozen unlabelled pilot support continued controller assessment, but independent source-linked closure evaluation is still required before claiming reliable autonomous research-gap triage.

\subsubsection*{AI assistance disclosure.}
Generative AI assisted writing, code development, and abstract-level inspection. Generated text and software tests are not treated as independent expert annotation or scientific evidence.

\begin{thebibliography}{17}
\bibitem{swanson1986} Swanson, D.R.: Fish oil, Raynaud's syndrome, and undiscovered public knowledge. Perspectives in Biology and Medicine 30(1), 7--18 (1986)
\bibitem{baek2025} Baek, J., Jauhar, S.K., Cucerzan, S., Hwang, S.J.: ResearchAgent: Iterative research idea generation over scientific literature with large language models. In: NAACL, pp. 6709--6738 (2025). \url{https://doi.org/10.18653/v1/2025.naacl-long.342}
\bibitem{mullerbloch2015} M\"uller-Bloch, C., Kranz, J.: A framework for rigorously identifying research gaps in qualitative literature reviews. In: ICIS, Research Methods, paper 2 (2015). \url{https://aisel.aisnet.org/icis2015/proceedings/ResearchMethods/2/}
\bibitem{wadden2020} Wadden, D., et al.: Fact or fiction: Verifying scientific claims. In: EMNLP, pp. 7534--7550 (2020). \url{https://doi.org/10.18653/v1/2020.emnlp-main.609}
\bibitem{wadden2022} Wadden, D., et al.: SciFact-Open: Towards open-domain scientific claim verification. In: Findings of EMNLP, pp. 4719--4734 (2022). \url{https://doi.org/10.18653/v1/2022.findings-emnlp.347}
\bibitem{chan2024} Chan, C.S.J., et al.: Overview of the Context24 shared task on contextualizing scientific claims. In: Workshop on Scholarly Document Processing, pp. 12--21 (2024). \url{https://aclanthology.org/2024.sdp-1.3/}
\bibitem{contextdata} Chan, C.S.J., et al.: Contextualizing scientific claims: Context24 dataset. HuggingFace, revision \texttt{457d3b5cb4bb8ade34e37458f4900c6eae0959bb}. \url{https://huggingface.co/datasets/joelchan/contextualizing-scientific-claims} (accessed 1 October 2026)
\bibitem{robertson2009} Robertson, S., Zaragoza, H.: The probabilistic relevance framework: BM25 and beyond. Foundations and Trends in Information Retrieval 3(4), 333--389 (2009). \url{https://doi.org/10.1561/1500000019}
\bibitem{cohan2020} Cohan, A., et al.: SPECTER: Document-level representation learning using citation-informed transformers. In: ACL, pp. 2270--2282 (2020). \url{https://doi.org/10.18653/v1/2020.acl-main.207}
\bibitem{singh2023} Singh, A., et al.: SciRepEval: A multi-format benchmark for scientific document representations. In: EMNLP, pp. 5548--5566 (2023). \url{https://doi.org/10.18653/v1/2023.emnlp-main.338}
\bibitem{bordes2013} Bordes, A., et al.: Translating embeddings for modeling multi-relational data. In: Advances in Neural Information Processing Systems 26 (2013)
\bibitem{blondel2008} Blondel, V.D., et al.: Fast unfolding of communities in large networks. Journal of Statistical Mechanics, P10008 (2008). \url{https://doi.org/10.1088/1742-5468/2008/10/P10008}
\bibitem{geifman2017} Geifman, Y., El-Yaniv, R.: Selective classification for deep neural networks. In: Advances in Neural Information Processing Systems 30 (2017)
\bibitem{jayawardena2026} Jayawardena, D., Perera, K.: AI-driven cybersecurity frameworks for real-time intrusion detection and threat intelligence. Int. J. Comput. Sci. Inf. Syst. 11(7), 60--77 (2026). \url{https://doi.org/10.55640/ijcsis/Volume11Issue07-02}
\bibitem{rahaman2023} Rahaman, M.S., Islam, A., Cerny, T., Hutton, S.: Static-analysis-based solutions to security challenges in cloud-native systems: Systematic mapping study. Sensors 23(4), 1755 (2023). \url{https://doi.org/10.3390/s23041755}
\bibitem{olaya2024} Plazas Olaya, M.K., Vergara Tejada, J.A., Aedo Cobo, J.E.: Securing microservices-based IoT networks: Real-time anomaly detection using machine learning. J. Comput. Networks Commun., 9281529 (2024). \url{https://doi.org/10.1155/2024/9281529}
\bibitem{kampa2024} Kampa, S.: Navigating the landscape of Kubernetes security threats and challenges. J. Knowledge Learning Sci. Technol. 3(4), 274--281 (2024). \url{https://doi.org/10.60087/jklst.v3.n4.p274}
\end{thebibliography}
\end{document}
